% docs/nuclear_gradients.tex -- the NEVPT2 nuclear gradient, as implemented.
%
% Current-state, like every file in docs/: it describes what the tree does,
% and it carries NO measured numbers (gradients, residuals, condition numbers,
% walls, agreements with grad_num). Those live in performance.md,
% pc-nevpt2.md, reference-data.md and testing.md, and are cited here by file
% and section name. The constants that appear below (the difference-quotient
% steps, GRAD_NUM_TOL, PC_GAP_FLOOR, ...) are constants of the code, not
% measurements; why each has the value it has is cited, not restated.
%
% Build: pdflatex nuclear_gradients.tex (twice, for the cross-references).
% Needs amsmath, amssymb, booktabs, algorithm, algpseudocode (algorithmicx),
% hyperref, geometry, lmodern and microtype -- all in a standard TeX Live.
\documentclass[11pt]{article}

\usepackage[T1]{fontenc}
\usepackage{lmodern}
\usepackage{microtype}
% "--pc" must stay two hyphens in \texttt, not an en dash.
\DisableLigatures[-]{encoding = T1, family = tt*}
\usepackage[margin=2.5cm]{geometry}
\usepackage{amsmath,amssymb}
\usepackage{booktabs}
\usepackage{algorithm}
\usepackage[noend]{algpseudocode}
\usepackage[hidelinks]{hyperref}

\newcommand{\code}[1]{\texttt{#1}}
\newcommand{\G}[1]{\Gamma^{(#1)}}
\newcommand{\eps}{\varepsilon}
\newcommand{\Fca}{F^{\mathrm{ca}}}
\newcommand{\Fac}{F^{\mathrm{ac}}}
\newcommand{\heff}{h^{\mathrm{eff}}}
\newcommand{\Ecas}{E_{\mathrm{CAS}}}
\newcommand{\Ecorr}{E_{\mathrm{corr}}}
\newcommand{\cls}[1]{\ensuremath{\mathrm{#1}}}
\newcommand{\cons}[1]{\ensuremath{(\mathrm{#1})}}
\algrenewcommand\algorithmicrequire{\textbf{Input:}}
\algrenewcommand\algorithmicensure{\textbf{Output:}}

\title{The NEVPT2 nuclear gradient:\\
SC and PC, on CASCI and CASSCF orbitals}
\author{\code{nevpt2-gpu-demo}, \code{docs/nuclear\_gradients.tex}}
\date{}

\begin{document}
\sloppy
\maketitle

\begin{abstract}
The tree assembles the analytic nuclear gradient of the strongly-contracted
(SC) and, with \code{--pc}, the partially-contracted (PC) NEVPT2 energy, and
checks it against a numerical gradient of re-converged energies to
$10^{-6}\,E_h/a_0$ per component. The work is split on purpose. The GPU
(\code{nevpt2\_sc\_pseudodensity --dump-source-terms}) computes the two
\emph{source terms}, the energy's derivatives with respect to the orbital
rotations and the CI coefficients, in reverse mode, and writes them out.
Offline Python (\code{reference\_data/nevpt2\_grad.py}) reads them, solves
the Lagrangian's Z-vector equations, and takes the nuclear derivative of one
scalar along a frozen orthonormal path. A second offline script
(\code{nevpt2\_numgrad.py}) supplies the numerical gradient that the result is
gated against. This document covers the four combinations
$\{\text{SC},\text{PC}\}\times\{\text{CASCI},\text{CASSCF}\}$ and the
density-fitted SC case. For each it gives the equations the code evaluates
and pseudocode for the path through the tree. The derivation is
\code{docs/gradient-theory.md}, the GPU side in prose is
\code{docs/gradients.md}, and the energy itself is \code{docs/energy.tex}.
Like every file in \code{docs/}, this one states what the code \emph{does},
and it carries no measured numbers.
\end{abstract}

\tableofcontents

%======================================================================
\section{The four paths, and where they actually differ}
\label{sec:paths}

\begin{table}[htbp]
\centering
\begin{tabular}{@{}lll@{}}
\toprule
 & CASCI orbitals (RHF) & CASSCF orbitals \\
\midrule
SC & \S\ref{sec:sc-casci} (fitted: \S\ref{sec:df}) & \S\ref{sec:sc-casscf} \\
PC (\code{--pc}) & \S\ref{sec:pc-casci} & \S\ref{sec:pc-casscf} \\
\bottomrule
\end{tabular}
\caption{The four paths, and the section that walks through each.}
\label{tab:paths}
\end{table}

For the energy the two axes of Table~\ref{tab:paths} split the work at the
input and at the class energies (\code{docs/energy.tex},
\S1). For the gradient they split it differently:

\begin{itemize}
\item \textbf{CASCI vs.\ CASSCF changes the Lagrangian.} The energy is the same
  functional on both kinds of orbitals. What differs is how the orbitals were
  made, and that is what the orbital constraints state. On RHF orbitals the
  core/active/virtual partition holds because the RHF Fock matrix is
  block-diagonal over it (\S\ref{sec:casci}). On CASSCF orbitals it holds
  because the CASSCF orbital gradient vanishes (\S\ref{sec:casscf}). The
  first gives a CI equation decoupled from the orbital one. The second gives
  one coupled CP-CASSCF system. The golden's header field
  \code{reference\_kind} records which branch applies, and
  \code{nevpt2\_grad.py --reference} must name the same one.
\item \textbf{SC vs.\ PC changes one scalar.} The Lagrangian is a statement
  about the reference, not about the contraction scheme. \code{--pc}
  replaces $\Ecorr$, PySCF's SC sum or the GPU's SC walk, by the PC energy
  (block2's offline, the GPU's PC walk on the card). Through $\Ecorr$ it also
  replaces the two source terms. Everything else is unchanged: the constraint
  set, the Jacobian, the frozen path and the solve order. PC adds its own
  difference-quotient steps and a conditioning guard at every displaced CI
  vector (\S\ref{sec:pc-casci}).
\end{itemize}

\begin{table}[htbp]
\centering
\small
\begin{tabular}{@{}lll@{}}
\toprule
stage & where & equations \\
\midrule
the source terms $Y_{xy}$, $y_I$ & GPU, reverse mode (\code{--dump-source-terms}) & \S\ref{sec:src-gpu} \\
 & \emph{or} offline, central differences of PySCF/block2 & \S\ref{sec:src-fd} \\
the constraint Jacobian, $\partial(Hc)/\partial\kappa$ & offline, one $\kappa$ sweep & \S\ref{sec:src-fd} \\
the Z-vector solve & offline only & \S\ref{sec:casci}, \S\ref{sec:casscf} \\
the skeleton derivative and $\Lambda\cdot\partial g/\partial\xi$ & offline, frozen-path differences of $\Phi$ & \S\ref{sec:frozen} \\
the numerical gradient \code{grad\_num} & offline, re-converged energies & \S\ref{sec:numgrad} \\
\bottomrule
\end{tabular}
\caption{Who computes what. Only the first row has a GPU path. Nothing in
\code{src/} or \code{apps/} differentiates with respect to a nuclear
coordinate (\code{docs/gradients.md}, \S1).}
\label{tab:who}
\end{table}

\textbf{The Z-vector solve is offline by design.} The repo's shape is
``offline Python generates the reference, the C++ GPU demo reproduces it'',
and a PySCF solve inside the runtime would invert it. The GPU's job is the one
expensive derivative of the NEVPT2 energy, its source terms. A gradient built
from the GPU's dump \emph{is} the NEVPT2 nuclear gradient, and is gated as one
(\code{docs/testing.md}, ``The GPU as the source-term provider'').

%======================================================================
\section{Notation and conventions}
\label{sec:notation}

\paragraph{Orbital spaces.} $i,j$ label core orbitals, $t,u,v,w$ active
orbitals and $r,s$ virtual orbitals. In this document $p,q,x,y$ label a
\emph{general} MO, any of the three; \code{docs/energy.tex} uses $p,q$ for
active ones. The counts are $n_c$, $n_a$, $n_v$, $n_{\mathrm{mo}} = n_c + n_a
+ n_v$. On the CASCI branch a second partition is in play, the RHF occupied
space $O$ and virtual space $V$. The active space straddles it:
$\mathrm{act}_O = \mathrm{act}\cap O$ and $\mathrm{act}_V = \mathrm{act}\cap
V$ (\code{docs/gradient-theory.md}, \S0). $I,J$ label determinants of the CI
vector $c$. $\xi$ is one Cartesian nuclear coordinate.

\paragraph{Orbital rotations.} $C(\kappa) = C\,e^{\kappa}$ with $\kappa$ real
antisymmetric, so that
\begin{align}
  \delta h_{pq} &= \sum_x \bigl(\kappa_{xp} h_{xq} + \kappa_{xq} h_{px}\bigr),
  \label{eq:dh}\\
  \delta(pq|rs) &= \sum_x \bigl(\kappa_{xp}(xq|rs) + \kappa_{xq}(px|rs)
     + \kappa_{xr}(pq|xs) + \kappa_{xs}(pq|rx)\bigr),
  \label{eq:dg}
\end{align}
in chemists' notation. The offline code steps a single pair $(p,q)$, $p<q$,
by the exact Givens rotation $C\exp\bigl(\theta(E_{pq}-E_{qp})\bigr)$
(\code{givens}).

\paragraph{Fock matrices.} $F$ is the \textbf{RHF} Fock matrix, built from the
RHF density $2\,C P_O C^{\mathsf T}$ and so from $O$ alone. $f$ is the
\textbf{generalised} Fock matrix, built from the zeroth-order density:
\begin{equation}
  \heff_{pq} = h_{pq} + \sum_j \bigl[2(pq|jj) - (pj|jq)\bigr],\qquad
  f_{pq} = \heff_{pq} + \sum_{tu}\G1_{tu}\bigl[(pq|tu) - \tfrac12(pt|uq)\bigr],
  \qquad \eps_p = f_{pp}.
  \label{eq:fock}
\end{equation}
$\eps_i$ and $\eps_r$ are the golden's \code{e\_core} and \code{e\_virt}, and
they move with the CI vector through $\G1$.

\paragraph{Density matrices.} $\G1,\G2,\G3,\Fca,\Fac$ are those of
\code{docs/energy.tex}, \S2. $E_{tu} = \sum_\sigma a^\dagger_{t\sigma}
a_{u\sigma}$. $P = 1 - |c\rangle\langle c|$ projects the state out of the
determinant space.

\paragraph{Constants of the code.} $\code{GRAD\_NUM\_TOL} = 10^{-6}$ (the gate
against \code{grad\_num}, and every ladder's tolerance);
$\code{PC\_GAP\_FLOOR} = 10^{-6}\,E_h$ and $\code{PC\_ZERO\_TOL} = 10^{-8}$
(the degenerate-pair rule, \S\ref{sec:casci}); the difference-quotient steps
of Table~\ref{tab:steps}; and the ladder rule's window $(3.2, 4.8)$ and run
length $3$ (\S\ref{sec:checks}). All are in \code{nevpt2\_grad.py} and
\code{nevpt2\_numgrad.py}. Why each step has its value is
\code{docs/testing.md}, ``The step constants, measured''. Why the gate is not
this repo's usual $10^{-7}$ is \code{docs/testing.md}, ``The analytic gradient
against the numerical one'': both sides are finite differences.

\begin{table}[htbp]
\centering
\small
\begin{tabular}{@{}llll@{}}
\toprule
step & differences & SC & PC \\
\midrule
\code{KAPPA\_STEP} & $\Ecorr$, $F$, $f$, $Hc$ along one $\kappa$ pair & $2^{-17}$ & $10^{-4}$ \\
\code{CI\_STEP} & $\Ecorr$ along $Pe_I$ & $2^{-17}$ & $10^{-7}$ \\
\code{NUC\_STEP} & $\Phi$ along $\xi$, frozen path (bohr) & $2^{-13}$ & $0.0025$ \\
\code{STEP\_BOHR} & re-converged $E$ along $\xi$ (\code{nevpt2\_numgrad.py}) & $0.0025$ & $0.0025$ \\
\bottomrule
\end{tabular}
\caption{The difference-quotient steps. PC's CI step is not chosen for
truncation error: it is the step at which the PC metric's lifted null modes
stay below the cut (\S\ref{sec:pc-casci}).}
\label{tab:steps}
\end{table}

%======================================================================
\section{What is differentiated}
\label{sec:functional}

The total energy is $E = \Ecas + \Ecorr$, a function of the geometry, the MO
coefficients $C$ and the CI vector $c$ (Park Eq.~21). The amplitude
dependence drops out because the Hylleraas functional is stationary in the
amplitudes, for either reference. $\Ecas = \langle c|H|c\rangle$ includes the
nuclear repulsion and the core energy.

\paragraph{The SC finish, and its three multipliers.} Each SC class reaches
$C$ and $c$ only through three scalars per external tuple $t$: the norm $N_t$,
the active part $H_t$ of the Dyall expectation value, and the orbital-energy
sum $\Delta_t$ (\code{docs/energy.tex}, \S6). With the amplitude $T_t$,
\begin{equation}
\begin{gathered}
  T_t = -\frac{1}{\Delta_t + H_t/N_t},\qquad E_t = T_t N_t,\\
  P_t = \frac{\partial E}{\partial N_t} = 2T_t + \Delta_t T_t^2,\qquad
  Q_t = \frac{\partial E}{\partial H_t} = T_t^2,\qquad
  R_t = \frac{\partial E}{\partial \Delta_t} = N_t T_t^2,
\end{gathered}
  \label{eq:pqr}
\end{equation}
which are Park's Eqs.~36--38 and satisfy $P_tN_t + Q_tH_t = E_t$ (Eq.~40).
This is the only hand-differentiated expression on the SC side; everything
downstream is $P$, $Q$ and $R$ contracted back against integrals and densities
(\code{nevpt2.sc\_amplitudes}, \code{nevpt2.energy\_finish}).

\paragraph{The PC finish.} A PC class solves in the span of its internal
functions, $E_t = -c_\kappa\,b_t^{\mathsf T} A_t^{+} b_t$ with $A_t = K + \Delta_t S$ and
$b_t = S x_t + S_1 h_t$ (\code{docs/energy.tex}, Eq.~(29); $c_\kappa$ the
class's coupling factor). With $x = A_t^{+} b_t$, the solve the energy
already ran,
\begin{equation}
  \dot E_t = -c_\kappa\,\bigl(2\,x^{\mathsf T}\dot b_t - x^{\mathsf T}\dot A_t\,x\bigr),
  \label{eq:pcdot}
\end{equation}
so no eigenvalue or eigenvector derivative appears. This holds only where the
kept rank of $S$ is locally constant. The energy's gap refusal enforces that,
and at an eigenstate $K$ annihilates $S$'s dropped space, which makes the two
terms of Eq.~\eqref{eq:pcdot} exact (\code{docs/gradients.md}, \S7).

\paragraph{What carries the CI vector.} The RDMs, and through $\G1$ in
Eq.~\eqref{eq:fock} the orbital energies. $\eps_p$ is affine in $\G1$:
\begin{equation}
  \eps_p = \heff_{pp} + \sum_{tu} \G1_{tu}\,W_p[t,u],\qquad
  W_p[t,u] = (pp|tu) - \tfrac12(pt|up),\qquad p\notin\mathrm{act}.
  \label{eq:epsW}
\end{equation}
Holding the orbital energies fixed differentiates a different functional, and
that omission is invisible to any check whose two sides share it
(\code{docs/gradients.md}, \S5.5).

%======================================================================
\section{\texorpdfstring{Shared stage A: the source terms $Y$ and $y$}{Shared stage A: the source terms Y and y}}
\label{sec:src}

\subsection{What the Lagrangian reads}
Park's Eqs.~26--27 define the two source terms as derivatives of $\Ecorr$ at
\emph{fixed} constraints, with no solve involved:
\begin{equation}
  Y_{xy} = \frac{\partial \Ecorr}{\partial \kappa_{xy}},\qquad
  y_I = \frac{\partial \Ecorr}{\partial c_I}.
  \label{eq:YY}
\end{equation}
$\kappa$ is antisymmetric and $c$ is normalised, so the solve reads only two
combinations:
\begin{equation}
  A_{pq} = Y_{pq} - Y_{qp},\qquad Py = y - (c\cdot y)\,c.
  \label{eq:Apy}
\end{equation}
Every block of $A$ has a reader. Core--core and virtual--virtual feed
\cons{O2}'s multipliers, active--active feeds the gauge identity, and the
three mixed blocks feed the orbital Z-vector equation
(\code{docs/gradient-theory.md}, \S2.1). The source terms are the same on both
branches: they differentiate the NEVPT2 energy, which does not know how its
orbitals were made. There are two providers.

\subsection{Provider 1: the GPU, by reverse mode}
\label{sec:src-gpu}

\paragraph{The CI source term.} \code{scCiSourceDevice} runs the class
energies' own slab walk once more. Each class is a list of norm and $H$ terms
in the same subscript strings as its energy einsums. One routine evaluates the
list and differentiates it by rewriting the subscripts per operand, so no
derivative is hand-transcribed. The walk leaves an accumulator $M_G =
\partial E/\partial G$ per metric and Dyall block $G$. Every block is linear
in the RDMs, so a reverse-topological sweep through the blocks' own builders
gives five coefficient tensors:
\begin{equation}
  \bar\rho = \frac{\partial E}{\partial\rho} = \sum_G
    \Bigl(\frac{\partial G}{\partial\rho}\Bigr)^{\mathsf T} M_G,
  \qquad \rho\in\{\G1,\G2,\G3,\Fca,\Fac\}.
  \label{eq:rdmadj}
\end{equation}
The orbital energies add one more $\G1$ coefficient. Each class scatters its
$R_t$ onto the $\eps_p$ its $\Delta_t$ reads, and the sum enters through
Eq.~\eqref{eq:epsW}:
\begin{equation}
  \bar\Gamma^{(1)}_{tu} \mathrel{+}= \sum_{p\notin\mathrm{act}}
    \bar\eps_p\,W_p[t,u],\qquad
  \bar\eps_p = \sum_t R_t\,\frac{\partial\Delta_t}{\partial\eps_p}.
  \label{eq:epsadj}
\end{equation}
Then
\begin{equation}
  y_I = \sum_\rho \Bigl\langle \bar\rho,\ \frac{\partial\rho}{\partial c_I}\Bigr\rangle,
  \label{eq:ysigma}
\end{equation}
taken without ever forming $\partial\G3/\partial c$ (an $n_a^6\,n_{\mathrm{det}}$
object) or the 4-RDM. One 3-body sigma vector carries $\bar\Gamma^{(1)}$,
$\bar\Gamma^{(2)}$ and $\bar\Gamma^{(3)}$, because the trace of a one-body
generator is the electron count on any determinant. A forward and a dagger
4-body contracted sigma for each $F$ order carry the digests, through one
$(n_a^2, n_{\mathrm{det}})$ intermediate (\code{docs/gradients.md}, \S5.3--5.5;
\code{nevpt2.sigma\_vector}).

\paragraph{The orbital source term.} Along a rotation at fixed $c$, every
stored array with an index the rotation touches moves. The stored two-electron
blocks move by Eq.~\eqref{eq:dg}, slot by slot. The one-electron blocks move,
their core dressing promoted with them. The orbital energies move. And the
digests move, because they are the active integrals contracted with the
4-RDM. The pairing is
\begin{equation}
  \langle\kappa, Y\rangle =
    \underbrace{\langle \bar I,\ \dot I(\kappa)\rangle
    + \sum_t R_t\,\dot\Delta_t(\kappa)}_{\text{frozen digests}}
    + \underbrace{\langle \bar v,\ \dot v(\kappa)\rangle}_{\text{the digests' term}},
  \label{eq:Ypair}
\end{equation}
where $I$ is every stored block, $\bar I$ its pseudodensity (Park Eq.~41 and
the Eqs.~42--47 intermediates), and $\bar v = \partial E/\partial v$ through
$\Fca,\Fac$ (with $v$ the active block, \code{docs/energy.tex},
Eq.~(1)). $\dot I(\kappa)$ is the \emph{promotion rule}: replace one slot of a
block by a general index and contract it with $\kappa$. Each right-hand side is
one of the gradient sidecar's one-general-index arrays
(Appendix~\ref{app:inventory}). $\bar v$ comes from running the digest build's
$f3$ path backwards, the \emph{transpose} (\code{buildDigestAdjointDevice}).
The plain build's \code{produce} kernel regenerates $L$ tile by tile, so the
4-RDM is never formed (\code{docs/gradients.md}, \S6). The active--active part
of the directly contractible terms is symmetric, so it is invisible to an
antisymmetric $\kappa$, and the matrix $Y$ is assembled by probing
Eq.~\eqref{eq:Ypair} along basis directions (\code{scFullRotationMatrix}).

\paragraph{PC.} The same walk with each PC class's solve differentiated
through its own $S$, $K$ and $S_1$, Eq.~\eqref{eq:pcdot} read as an adjoint:
\begin{equation}
  \frac{\partial E}{\partial K} = U\,\mathrm{sym}(m)\,U^{\mathsf T},\quad
  \frac{\partial E}{\partial S} = \mathrm{sym}\bigl(U\,\bar T^{\mathsf T}\bigr)
     + U\,\mathrm{sym}\Bigl(\textstyle\sum_t \Delta_t m_t\Bigr)U^{\mathsf T},\quad
  \frac{\partial E}{\partial S_1} = U\,\bar t_1^{\mathsf T},\quad
  R_t = \frac{\partial E}{\partial\Delta_t},
  \label{eq:pcadj}
\end{equation}
with $U$, $T$, $t_1$ the solve's eigenbasis quantities (\code{docs/energy.tex},
Eqs.~(30)--(33)), $m_t$ the tuple's weight in that basis, and $\bar T$,
$\bar t_1$ the slab einsum's adjoint in its basis operand. $K$ is SC's own
Dyall blocks and $S$, $S_1$ are permutations of SC's metrics, so
Eq.~\eqref{eq:rdmadj}'s sweep carries them down unchanged. A class that the
energy's gap or denominator test refuses carries that refusal out of every
entry point; it never contributes a zero (\code{docs/gradients.md}, \S5.6,
\S6.1).

\paragraph{The gates the GPU holds itself to.} Both source terms are also
computed in forward mode, a tangent $\dot E(u)$ along a CI direction and
$\dot E(\kappa)$ along a rotation. The two modes share no code past the slab
walk, so their agreement is an exact identity with no step size in it:
\begin{equation}
  \langle u, y\rangle = \dot E(u),\qquad
  \langle\kappa, Y\rangle = \dot E_{\mathrm{frozen}}(\kappa) + \dot E_{\mathrm{digest}}(\kappa).
  \label{eq:duality}
\end{equation}
The dualities cannot see a term both modes omit. The orbital energies'
response and the digests' term are therefore also held to PySCF's exact
derivative and to finite differences of the unmodified energy
(\code{docs/gradients.md}, \S2 and \S5.5).

\begin{algorithm}[htbp]
\caption{\code{nevpt2\_sc\_pseudodensity --dump-source-terms} \textit{P} (\code{apps/sc\_pseudodensity/source\_dump.cppm})}
\label{alg:dump}
\begin{algorithmic}[1]
\Require golden $g$ and its \code{.gradgold} sidecar; an exact digest (\code{--cublas}, \code{--ozaki}, \code{--fused-digest} refused); \code{--pc} optional
\State $(\G3,\Fca,\Fac) \gets$ \code{buildRdmsDevice} \Comment{\code{docs/energy.tex}, Alg.~2}
\State $W \gets$ \code{orbitalEnergyResponse}(sidecar) \Comment{Eq.~\eqref{eq:epsW}}
\State $\bar v \gets$ the digests' transpose \Comment{\code{scDigestAdjointDevice}, or PC's}
\State $Y \gets$ probe Eq.~\eqref{eq:Ypair} along every MO pair \Comment{\code{scFullRotationMatrix}}
\State $y \gets$ Eqs.~\eqref{eq:rdmadj}--\eqref{eq:ysigma} with $W$ \Comment{\code{scCiSourceDevice}, or \code{pcCiSourceDevice}}
\State along six seeded directions $\kappa_d$: $\dot E_{\mathrm{frozen}}(\kappa_d)$, $\dot E_{\mathrm{digest}}(\kappa_d)$
\State \textbf{gate:} the matrix is the pairing; $\dot E(\kappa_d) = \langle\kappa_d, Y\rangle$ to $10^{-7}$; $y$'s walk reproduces $\Ecorr$ to $10^{-7}$; PC: the integral halves' duality too
\State write $P$ (\code{.gradsrc}, an \code{.npz}): \code{kappa\_source}, \code{ci\_source}, \code{ci}, $n_c,n_a,n_v$, \code{e\_corr}, \code{pc}, \code{edot\_*}; read it back bit for bit
\end{algorithmic}
\end{algorithm}

\paragraph{Carrying the dump into the gradient's basis.} The GPU knows one
active basis, the golden's natural orbitals. On the CASCI branch
\code{nevpt2\_grad.py} holds the state in the eigenbasis of $F$'s active block
(\S\ref{sec:casci}), reached by an active rotation $W_a$ that carries $c$
along. \code{load\_source\_terms} therefore first checks that the dump's
\code{ci} is the golden's up to a sign $\sigma$ (to \code{DUMP\_CI\_TOL}),
that its partition is the state's, that its \code{e\_corr} matches, and that
its scheme (SC or PC) is the run's. It then transforms
\begin{equation}
  Y' = W^{\mathsf T} Y\,W,\quad W = 1\oplus W_a,\qquad
  y' = \mathcal T(W_a)\,(\sigma y),
  \label{eq:dumprot}
\end{equation}
with $\mathcal T(W_a)$ the CI transform for the orbital rotation (orthogonal,
so $y$ transforms as $c$ does). $E$ is even in $c$, so $y$ is odd, and the sign
goes on $y$. On the CASSCF branch the state is the golden's own basis and
$W_a = 1$.

\subsection{Provider 2: offline central differences}
\label{sec:src-fd}

\code{--source pyscf}, the default, differentiates PySCF's SC energy, or
block2's PC energy with \code{--pc}, and needs no card. One pass over every
non-gauge pair $k = (p,q)$ (\code{kappa\_sweep}) produces three objects at
once:
\begin{align}
  A_{pq} &= \frac{E(C\,G_{pq}(h)) - E(C\,G_{pq}(-h))}{2h},
  \label{eq:Afd}\\
  M_{km} &= \frac{\partial G_m}{\partial\kappa_k},\qquad
  Q_{kI} = \frac{\partial (Hc)_I}{\partial\kappa_k},
  \label{eq:MQ}
\end{align}
at $h = \code{KAPPA\_STEP}$ and fixed $c$, with $G_{pq}$ the Givens rotation
and $G_m$ the orbital constraints of \S\ref{sec:casci}--\ref{sec:casscf}. $M$
is the constraint Jacobian. $Q$ makes the CI relaxation $Qz$ available for any
$z$ found later. $M$ is differenced rather than written out because $f$ is not
covariant under a rotation that mixes the core and active index ranges. The
CI source term is differenced along directions \emph{orthogonal to the state}
and inside the state's spatial irrep (\code{allowed\_determinants}):
\begin{equation}
  y_I = \|Pe_I\|\;\frac{\Ecorr(c + h\,d_I) - \Ecorr(c - h\,d_I)}{2h},\qquad
  d_I = \frac{Pe_I}{\|Pe_I\|},
  \label{eq:yfd}
\end{equation}
at $h = \code{CI\_STEP}$, with $\eps$ rebuilt from Eq.~\eqref{eq:fock} at
every displaced $c$. Both restrictions are load-bearing. Off the unit sphere
the hole-side metrics' constants are no longer hole RDMs and the SC finish has
poles. Outside the irrep a symmetry-empty perturber fills linearly in $h$, and
the limit, though it converges, is not linear in the direction
(\code{nevpt2\_grad.py}, \code{pyscf\_ci\_source}). When a dump supplies $Y$,
the sweep still runs to build $M$ and $Q$, but with $\Ecorr$ dropped, and the
dump's $A$ is added to the sweep's $\Ecas$ half.

\code{--source both} solves twice on one sweep, once from the dump and once
from PySCF's or block2's quotients. It compares the two source terms block by
block and prints both gradients beside \code{grad\_num}, so a difference
between them can only come from the source terms.

%======================================================================
\section{\texorpdfstring{Shared stage B: the frozen orthonormal path, and $\Phi$}{Shared stage B: the frozen orthonormal path, and Phi}}
\label{sec:frozen}

Given the multipliers $\Lambda$, the gradient is the nuclear derivative of one
scalar along a path on which the wavefunction parameters do not respond:
\begin{equation}
  C(\xi) = S(\xi)^{-1/2}\,S(\xi_0)^{1/2}\,C(\xi_0),\qquad c(\xi) = c(\xi_0),
  \label{eq:frozen}
\end{equation}
with $S$ the AO overlap (\code{frozen\_coeff}). Then $C(\xi)^{\mathsf
T}S(\xi)C(\xi) = 1$ at every geometry. The orthonormality constraint holds
identically along the path and needs no multiplier, so no energy-weighted
density term appears. With
\begin{equation}
  \Phi(\xi) = \Ecas(\xi) + \Ecorr(\xi) + \sum_m \lambda_m G_m(\xi) + \Lambda_{\mathrm{CI}}(\xi),
  \label{eq:phi}
\end{equation}
where $\Lambda_{\mathrm{CI}}$ is the CI constraint's term ($\langle z|H|c\rangle$
on the CASCI branch, $z\cdot g_{\mathrm{ci}}$ on the CASSCF one), every
quantity is evaluated at $(C(\xi), c)$ on the displaced molecule's own
integrals. The multipliers satisfy their equations, so
\begin{equation}
  \frac{dE}{d\xi} = \frac{d\Phi}{d\xi}
  \approx \frac{4\,g(h/2) - g(h)}{3},\qquad
  g(h) = \frac{\Phi(\xi + h) - \Phi(\xi - h)}{2h},
  \label{eq:dphi}
\end{equation}
at $h = \code{NUC\_STEP}$. The Richardson bar $|R - g(h/2)|$ is printed with
the gradient. The two halves are reported separately:
\code{get\_ham\_response} is the derivative of $\Ecas + \Ecorr$, and
\code{get\_LdotJnuc} is the derivative of the constraint terms.

\textbf{This is a Lagrangian gradient whose skeleton derivatives are
differenced, not contracted against derivative integrals}, and that is a stated
limitation. All the analytic content is exercised: the constraint set, the
Z-vector equations and the stationarity that makes a frozen path sufficient.
Replacing the differences with PySCF's derivative integrals would need the
NEVPT2 relaxed densities, and no flag writes those out (\S\ref{sec:limits}).
Nothing re-solves along the path, so every Cartesian component is
differenced, including the ones that are zero by symmetry. They come out as
the zeros they are, which is a free check (\S\ref{sec:numgrad}).

\begin{algorithm}[htbp]
\caption{\code{Gradients.kernel}: the frozen-path derivative}
\label{alg:kernel}
\begin{algorithmic}[1]
\Require the state $(C_0, c)$ at $\xi_0$; multipliers $\lambda$, $z$; step $h$
\State $S_0^{1/2} \gets$ \code{sym\_power}$(S(\xi_0), \tfrac12)$
\For{each atom and Cartesian component $\xi$} \Comment{all $3N$, zero-by-symmetry included}
  \For{$s \in \{h, h/2\}$}
    \For{$\pm$}
      \State build the molecule at $\xi \pm s$; its RHF object, integrals and, fitted, its own \code{with\_df}
      \State $C \gets S(\xi\pm s)^{-1/2} S_0^{1/2} C_0$ \Comment{Eq.~\eqref{eq:frozen}}
      \State $\Phi_\pm \gets \Ecas + \Ecorr + \lambda\cdot G + \Lambda_{\mathrm{CI}}$ at $(C, c)$ \Comment{Eq.~\eqref{eq:phi}}
    \EndFor
    \State $g(s) \gets (\Phi_+ - \Phi_-)/2s$
  \EndFor
  \State $\partial E/\partial\xi \gets [4g(h/2) - g(h)]/3$;\ bar $\gets |{\cdot} - g(h/2)|$
\EndFor
\end{algorithmic}
\end{algorithm}

%======================================================================
\section{The CASCI-branch Lagrangian and its solve}
\label{sec:casci}

\subsection{The constraints}
On RHF orbitals, canonicalised with \code{cas\_natorb=True}, core, active and
virtual are each spanned by a fixed set of RHF canonical orbitals
(\code{docs/gradient-theory.md}, \S1.2). The constraint set is
\begin{align}
  \cons{O1}&:\ F_{pq} = 0\ \text{for $p$, $q$ in different spaces of
     $\{\text{core},\mathrm{act}_O,\mathrm{act}_V,\text{virt}\}$},\\
  \cons{O2}&:\ f_{ij} = 0\ (i\ne j\ \text{core}),\quad f_{rs} = 0\ (r\ne s\
     \text{virtual}),\\
  \cons{C1}&:\ (H - \Ecas)\,c = 0,\quad c^{\mathsf T}c = 1,
\end{align}
and the Lagrangian is
\begin{equation}
  L = \Ecas + \Ecorr + \sum_{m\in\cons{O1}} \lambda_m F_{p_mq_m}
    + \sum_{m\in\cons{O2}} \lambda_m f_{p_mq_m}
    + \langle z|H - \Ecas|c\rangle.
  \label{eq:Lcasci}
\end{equation}
\cons{O1}'s blocks between $O$ and $V$ are RHF Brillouin conditions. Its blocks
within $O$ (core--$\mathrm{act}_O$) and within $V$ ($\mathrm{act}_V$--virtual)
hold by the active-space \emph{selection} (\S\ref{sec:limits}).

\paragraph{The gauges are $\mathrm{act}_O$-internal and $\mathrm{act}_V$-internal rotations, and no others.}
An $\mathrm{act}_O$--$\mathrm{act}_V$ rotation leaves the energy alone, being an
active--active rotation of a full CI, but it moves the RHF occupied space and
therefore $F$. So $F_{\mathrm{act}_O,\mathrm{act}_V} = 0$ is a genuine
Brillouin condition with a multiplier of its own. For that reason the state is
held in the basis that diagonalises $F$'s active block, with the CI vector
transformed to match, so that $\mathrm{act}_O$ and $\mathrm{act}_V$ are index
ranges (\code{build\_state}; \code{docs/gradient-theory.md}, \S2.6). The
rebuild is checked: $\Ecorr$ must not move under the rotation, and $F$'s
active block must come out diagonal. \code{kappa\_pairs} is every $p<q$ except
those two gauge blocks. \code{constraints} pairs them one to one with \cons{O1}
(between spaces) and \cons{O2} (inside core or inside virtual), and asserts
the count, so stationarity is a \emph{square} system.

\subsection{The stationarity equations}
Differentiating Eq.~\eqref{eq:Lcasci}:
\begin{align}
  \frac{\partial L}{\partial c}&:\quad P(H - \Ecas)P\,z = -P\bigl(y + y^c\bigr),
  \label{eq:ciz}\\
  \frac{\partial L}{\partial \kappa_k}&:\quad A_k + (Qz)_k + \sum_m M_{km}\lambda_m = 0
  \quad\text{for every non-gauge pair $k$},
  \label{eq:orbz}
\end{align}
with $A$, $M$, $Q$ from Eqs.~\eqref{eq:Afd}--\eqref{eq:MQ} (or $A$ from the
dump) and $y$ masked to the state's irrep. $F$ is built from the RHF density
and has no CI dependence, so Eq.~\eqref{eq:ciz} contains no $\cons{O1}$ term.
That is the decoupling of \code{docs/gradient-theory.md}, \S2.2. One coupling
survives. $f$ is built from $\G1$ and so does depend on $c$, which puts
\cons{O2}'s multipliers into the CI right-hand side as a one-body sigma vector
(\code{ci\_coupling}):
\begin{equation}
  y^c_I = \sum_{tu}\bigl(\mathcal W + \mathcal W^{\mathsf T}\bigr)_{tu}\,(E_{tu}c)_I,\qquad
  \mathcal W_{tu} = \sum_{m\in\cons{O2}} \lambda_m\bigl[(p_mq_m|tu) - \tfrac12(p_mt|uq_m)\bigr].
  \label{eq:yc}
\end{equation}

\paragraph{The solve.} Eq.~\eqref{eq:ciz} goes to preconditioned conjugate
gradient: $P(H-\Ecas)P$ is positive semidefinite on the projected space. The
preconditioner projects, and the answer is projected again, $z \gets z -
(c\cdot z)c$. $|c\rangle$ is an exact null vector, so a Jacobi preconditioner
would otherwise let $z$ drift along it at no cost to the residual. $\Phi$ drops
$-\Ecas'(\xi)\,(z\cdot c)$ on the assumption that $z\cdot c = 0$. The orbital
system, Eq.~\eqref{eq:orbz}, is square and indefinite by construction: the
\cons{O2} diagonal $\varepsilon_q - \varepsilon_p$ changes sign with the pair
order. It gets a direct least-squares solve (\code{JACOBIAN\_RCOND}), with
GMRES through the same operator beside it as an independent residual. The two
equations are iterated to a fixed point, because $y^c$ reads $\lambda$ and the
orbital right-hand side reads $Qz$. The iteration count is printed.

\paragraph{Degenerate \cons{O2} pairs.} Where $|\eps_p - \eps_q| <
\code{PC\_GAP\_FLOOR}$, the pair's multiplier is a division the state does not
determine. If $|A_{pq}| \le \code{PC\_ZERO\_TOL}$, the pair and its constraint
are dropped together, which keeps the system square, and the multiplier is
left at zero. An abelian subgroup puts two degenerate partners in different
irreps, so the source term vanishes there too and zero is the right
representative. Otherwise the run is refused rather than divided through
(\code{degenerate\_constraints}).

\paragraph{The theory's triangular order falls out.} On \cons{O2}'s rows $M$
is diagonal in the pseudocanonical gauge. Those rows reduce to Park's
Eqs.~56--57 (Appendix~\ref{app:z}), the selection-boundary rows to two small
Sylvester equations, and the rest to CPHF on the $O$--$V$ block
(\code{docs/gradient-theory.md}, \S2.3--2.5). The code imposes none of that.
It solves the one square system, so a response block that the triangular form
assumes to vanish is never relied on.

\begin{algorithm}[htbp]
\caption{\code{Gradients} (CASCI branch): \code{prepare}, \code{solve\_multipliers}, \code{kernel}}
\label{alg:casci}
\begin{algorithmic}[1]
\Require the rebuilt state; source terms $(A, y)$ from a dump or \code{None}
\State rotate the active block to $F$'s eigenbasis, CI vector along; check $\Ecorr$ unmoved and \cons{O1}, \cons{O2} \Comment{\S\ref{sec:casci}}
\State $(A_{\mathrm{sweep}}, M, Q) \gets$ \code{kappa\_sweep} \Comment{Eqs.~\eqref{eq:Afd}--\eqref{eq:MQ}; $\Ecorr$ dropped under a dump}
\State $A \gets A_{\mathrm{sweep}}$, or $A_{\mathrm{sweep}}^{\Ecas} + A_{\mathrm{dump}}$; \ $y \gets$ Eq.~\eqref{eq:yfd} or the dump's
\State $\lambda \gets 0$;\ $z \gets 0$
\Repeat \Comment{the CI/\cons{O2} fixed point}
  \State $b_{\mathrm{ci}} \gets P\bigl(y + y^c(\lambda)\bigr)$ \Comment{Eq.~\eqref{eq:yc}}
  \State $z \gets \textsc{Pcg}\bigl(P(H-\Ecas)P,\ -b_{\mathrm{ci}}\bigr)$;\ $z \gets Pz$
  \State drop degenerate \cons{O2} pairs, or refuse
  \State $\lambda \gets \textsc{Lstsq}\bigl(M,\ -(A + Qz)\bigr)$;\ GMRES beside it
\Until{the step in $(z,\lambda)$ is below $10^{-12}$}
\State report the gauge identity: $\max|A + Qz|$ over the $\mathrm{act}_O$-, $\mathrm{act}_V$-internal rotations
\State $\partial E/\partial\xi \gets$ Algorithm~\ref{alg:kernel} with $\lambda\cdot G = \lambda\cdot(F, f)$ and $\Lambda_{\mathrm{CI}} = z\cdot(H(\xi)c)$
\end{algorithmic}
\end{algorithm}

%======================================================================
\section{The CASSCF-branch Lagrangian and its solve}
\label{sec:casscf}

On converged CASSCF orbitals, Park's Eq.~23 holds as written
(\code{docs/gradient-theory.md}, \S1.4). Its first line, CASSCF
stationarity, replaces \cons{O1}:
\begin{equation}
  L = \Ecas + \Ecorr + Z\cdot g_{\mathrm{orb}}(\kappa, c) + z\cdot g_{\mathrm{ci}}(\kappa, c)
    + \sum_{m\in\cons{O2}} z^c_m\,f_{p_mq_m},
  \label{eq:Lcasscf}
\end{equation}
with $g = (g_{\mathrm{orb}}, g_{\mathrm{ci}})$ PySCF's second-order CASSCF
gradient, \code{newton\_casscf.gen\_g\_hop}'s \code{g\_all}. The orbital part
$g_{\mathrm{orb}}$ runs over the non-redundant core--active, core--virtual and
active--virtual rotations (constraint kind \cons{S1}), and $g_{\mathrm{ci}} =
2(H - \Ecas)c$ is \cons{C1}. There is no $F$, no \cons{O1}, no
$\mathrm{act}_O/\mathrm{act}_V$ split and no rotation of the state. Every
active--active rotation is a gauge.

\paragraph{The solve, in two steps and one pass.} First $z^c$, alone, from
\cons{O2}'s own rows. At a stationary point no other constraint responds to a
core--core or virtual--virtual rotation (\code{solve\_o2}, with the same
degenerate-pair rule as \S\ref{sec:casci}):
\begin{equation}
  M_{oo}\,z^c = -A_{oo}.
  \label{eq:zc}
\end{equation}
Then one coupled system (\code{solve\_coupled}), with $J = \partial g/\partial
x$ (\code{gen\_g\_hop}'s \code{h\_op}, the CASSCF orbital+CI Hessian in its own
step variables $x$), solved on a basis $B$: the identity on the orbital
variables, and an orthonormal basis of the in-sector determinants with $c$
projected out:
\begin{equation}
\begin{gathered}
  \bigl(B^{\mathsf T}JB\bigr)^{\mathsf T}\,\ell = -B^{\mathsf T}b,\qquad
  \begin{pmatrix} Z\\ z\end{pmatrix} = B\ell,\\
  b = \begin{pmatrix} -\bigl(A_{pq} + \sum_{o\in\cons{O2}} M_{ko}\,z^c_o\bigr)_{k=(p,q)}\\[2pt]
                      P\bigl(y + y^c(z^c)\bigr)\end{pmatrix}.
\end{gathered}
  \label{eq:cp}
\end{equation}
The minus sign on the orbital rows is the Givens convention: the rotation by
$\theta$ on $(p,q)$ is \code{gen\_g\_hop}'s step $x_{(q,p)} = -\theta$. By
default $B^{\mathsf T}JB$ is built densely, one \code{h\_op} product per
column, and LU-factorised. The solve is exact, with no preconditioner, level
shift or threshold. Its symmetric part is also diagonalised, and the modes
nearest singular are printed with their weight on the right-hand side, so a
near-null or negative mode is seen directly rather than inferred.
\code{--cp-solve cg} runs the base class's conjugate gradient through the same
operator, for a case too large to build densely.

\begin{algorithm}[htbp]
\caption{\code{CasscfGradients} (CASSCF branch)}
\label{alg:casscf}
\begin{algorithmic}[1]
\Require the state on \code{--mo-from}'s orbitals, unrotated; source terms as Algorithm~\ref{alg:casci}
\State $(A, M_{oo}) \gets$ \code{kappa\_sweep} \Comment{only \cons{O2}'s columns of $M$ are differenced}
\State $(g, J) \gets$ \code{gen\_g\_hop}$(C, c)$;\ report $\|g_{\mathrm{orb}}\|$, $\|g_{\mathrm{ci}}\|$ (the residual stationarity)
\State $z^c \gets$ Eq.~\eqref{eq:zc} \Comment{degenerate pairs left at zero, or refused}
\State $B \gets$ orbital variables $\oplus$ (in-sector determinants $\perp c$)
\State $J_{\mathrm{red}} \gets B^{\mathsf T}JB$, column by column;\ $\ell \gets J_{\mathrm{red}}^{-\mathsf T}(-B^{\mathsf T}b)$ \Comment{Eq.~\eqref{eq:cp}}
\State report $\|J_{\mathrm{red}}^{\mathsf T}\ell + B^{\mathsf T}b\|$, the same $\ell$ through \code{h\_op}, $J$'s asymmetry and its spectrum
\State $(Z, z) \gets B\ell$
\State $\partial E/\partial\xi \gets$ Algorithm~\ref{alg:kernel} with $\lambda\cdot G = Z\cdot g_{\mathrm{orb}} + z^c\cdot f$ and $\Lambda_{\mathrm{CI}} = z\cdot g_{\mathrm{ci}}$, $g$ from \code{gen\_g\_hop} at each displaced geometry
\end{algorithmic}
\end{algorithm}

%======================================================================
\section{The numerical gradient, and the gate}
\label{sec:numgrad}

\code{nevpt2\_numgrad.py} differences re-converged energies built by
\code{generate\_golden.\_build\_state} itself, at the same settings as the
golden. Those settings are single-threading, a $10^{-12}$ Davidson tolerance,
the point group detected, and the same irrep-selected active space. For each
probed component, with $h = \code{STEP\_BOHR}$,
\begin{equation}
\begin{gathered}
  g(s) = \frac{E(\xi+s) - E(\xi-s)}{2s},\qquad
  R(s) = \frac{4\,g(s/2) - g(s)}{3},\\
  \code{grad\_num} = R(h/2),\qquad
  \code{grad\_num\_residual} = |R(h/2) - R(h)|,
\end{gathered}
  \label{eq:numgrad}
\end{equation}
so each component costs six displaced energies. The residual is an error bar
on the stored number.

\paragraph{In-plane displacements only, and the refusal.} The active space is
selected by irrep, and irrep labels exist only while the point group does. A
displacement whose detected point group differs from the reference's is
refused (\code{\_check\_point\_group}). So is one whose core or active irrep
labels differ at an unchanged group (\code{\_check\_selection}). The
components such a displacement would probe are zero by symmetry at a planar
geometry. \code{--components auto} records them as exact zeros and leaves them
to the analytic side as a free check (\code{docs/gradient-theory.md}, \S4.3).

\paragraph{On a CASSCF reference} (\code{--reference casscf}) every displaced
state is re-optimised from a cold start through the generator's own optimiser,
then finished to a tighter orbital gradient, then given the golden's
\code{--mo-from} CASCI and NEVPT2. A displaced optimisation that lands on
another stationary point would give a number from a different reference, so
every displaced state passes a continuity guard before it is differenced. The
guard refuses an energy off its first-order prediction, an active subspace or
CI vector that moved too far, and an unconverged optimiser. A warm start on a
few geometries probes for a second solution (\code{docs/reference-data.md},
``The numerical gradient of a CASSCF reference''). Its PC twin runs no sweep:
\code{--pc-states} differences the GPU's PC energies on the finished sweep's
kept states (\code{pc\_on\_states.py}).

\paragraph{The gate} (\code{nevpt2\_grad.py}'s last step) compares each probed
component $\alpha$ with \code{grad\_num}:
\begin{equation}
  \max_\alpha \bigl|g_\alpha - \code{grad\_num}_\alpha\bigr| < \code{GRAD\_NUM\_TOL} = 10^{-6},
  \label{eq:gate}
\end{equation}
and prints \code{PASS:} or \code{FAIL:}, exiting 1 on a failure. The source is
the sidecar's \code{grad\_num}, or an \code{.npz} named with
\code{--grad-num}; the file must be of the same case, reference and functional
(\code{grad\_num}, \code{grad\_num\_pc}, or \code{grad\_casci} under
\code{--casci-only}). Two more checks are free and always reported: the
gradient summed over atoms (translational invariance), and $|g|$ on the axes
along which every atom shares a coordinate, i.e.\ zero by symmetry
(\code{zero\_axes}). \textbf{There is no ctest entry for the gate.} It would put
PySCF and block2 in the test path, and the gate needs no card. Only the GPU dump
it reads is a ctest entry (\code{docs/testing.md}, ``Why there is no ctest entry
for it'').

%======================================================================
\section{Path 1: SC-NEVPT2 on a CASCI reference}
\label{sec:sc-casci}

\paragraph{Offline, before.} The golden is the energy's
(\code{docs/energy.tex}, Path~1, conventional integrals here). Beside it,
\code{generate\_golden.py --gradgold} writes the \code{.gradgold} sidecar: the
one-general-index integral inventory of Appendix~\ref{app:inventory}, the
full-range one-electron and Fock matrices, $F$ and \code{mo\_coeff}.
\code{nevpt2\_numgrad.py} appends \code{grad\_num} to it, or writes an
\code{.npz} for a case whose committed sidecar carries none.

\begin{algorithm}[htbp]
\caption{Path 1, SC on CASCI orbitals}
\label{alg:path1}
\begin{algorithmic}[1]
\Statex \textit{Offline} (\code{generate\_golden.py --gradgold}; \code{nevpt2\_numgrad.py [--npz]})
\State golden and sidecar, as Path~1 of \code{docs/energy.tex} plus the inventory
\State \code{grad\_num} $\gets$ Eq.~\eqref{eq:numgrad} over the in-plane components
\Statex \textit{GPU} (\code{nevpt2\_sc\_pseudodensity --golden \textit{g} --tiles 3 --dump-source-terms \textit{P}})
\State $(Y, y) \gets$ Algorithm~\ref{alg:dump};\ \textbf{PASS} on its own gates
\Statex \textit{Offline} (\code{nevpt2\_grad.py --case \textit{c} --source dump|both --source-path \textit{P}})
\State rebuild the golden's state; refuse it unless $\Ecas$, $\Ecorr$ are the golden's to \code{REF\_STATE\_TOL}
\State carry $(Y, y)$ into the $F$-eigenbasis \Comment{Eq.~\eqref{eq:dumprot}}
\State $(\lambda, z) \gets$ Algorithm~\ref{alg:casci};\ $g \gets$ Algorithm~\ref{alg:kernel}
\State the residual-independent checks \Comment{\S\ref{sec:checks}}
\State \textbf{PASS} iff Eq.~\eqref{eq:gate} holds and every residual-independent check passes
\end{algorithmic}
\end{algorithm}

The equations are Eqs.~\eqref{eq:pqr}--\eqref{eq:dumprot},
\eqref{eq:frozen}--\eqref{eq:dphi}, \eqref{eq:Lcasci}--\eqref{eq:yc} and
\eqref{eq:numgrad}--\eqref{eq:gate}. Without \code{--source}, Algorithm~\ref{alg:path1}'s
GPU step is replaced by Eqs.~\eqref{eq:Afd} and \eqref{eq:yfd} and needs no
card. The cases are LiF CAS(4,6) and acrolein CAS(6,5), the planar $C_s$
gradient-validation pair (acrolein's sidecar is gitignored), and N2 CAS(4,4).
N2 is the deliberately symmetric case: six exactly degenerate inactive pairs,
and a linear molecule that keeps its point group under every displacement. Its
committed sidecar carries no \code{grad\_num}, so the route there is
\code{--npz}. The GPU dump's ctest entries are
\code{nevpt2\_sc\_pseudodensity\_dump\_\{n2\_ccpvdz\_cas44,
n2\_ccpvdz\_cas88, lif\_def2svp\_cas46\}}, plus the CAS(4,4) and LiF
\code{\_casscf} twins.
The verdicts are \code{docs/testing.md}, ``The analytic gradient against the
numerical one'' and ``The GPU as the source-term provider''.

\subsection{The density-fitted variant}
\label{sec:df}

\code{--df} differentiates the density-fitted SC energy on the \code{\_df}
golden's state: a fitted RHF and a \code{DFCASCI} sharing its \code{with\_df},
with \code{dfnevpt2.\_ERIS} called directly. Nothing else changes. The fitted
RHF is still an RHF, so \cons{O1}, \cons{O2}, \cons{C1}, the frozen path and
the solve are this section's. \textbf{The fit's own nuclear derivative is in the
quotient with no term written out}: each displaced geometry in
Algorithm~\ref{alg:kernel} builds its own fitting object, so the auxiliary
functions and the metric $(P|Q)$ move with the nuclei. An orbital rotation, by
contrast, never reaches the metric (\code{docs/gradient-theory.md}, \S3.8):
\begin{equation}
  (pq|rs) = \sum_L B^L_{pq}B^L_{rs},\qquad
  \delta B^L_{pq} = \sum_x \bigl(\kappa_{xp}B^L_{xq} + \kappa_{xq}B^L_{px}\bigr),
  \label{eq:dB}
\end{equation}
so the promotion rule acts on one $B$ factor at a time, and every inventory
array is a strip of $B$ against a block of $B$. The GPU's fitted dump
(\code{nevpt2\_sc\_pseudodensity --dump-source-terms} on the \code{\_df}
golden beside its \code{NEVPT2F1} sidecar) contracts that inventory out of the
whole $B$ and carries both terms of $Y$. The gate is the \code{NEVPT2F1}
sidecar's \code{grad\_num}, which \code{nevpt2\_numgrad.py --df} wrote. The
variant is SC, CASCI branch and N2 CAS(4,4) only. \code{--pc},
\code{--reference casscf}, \code{--casci-only} and \code{--grad-num} are
refused with it (\code{docs/testing.md}, ``The density-fitted gradient against
the numerical one'').

%======================================================================
\section{Path 2: PC-NEVPT2 on a CASCI reference}
\label{sec:pc-casci}

\paragraph{What \code{--pc} changes.} $\Ecorr$ becomes block2's PC energy
(\code{generate\_golden.\_block2\_icmr}) in the quotients and in $\Phi$. The
GPU's dump becomes the PC one (\code{--pc --dump-source-terms}, marked
\code{pc}; a dump of the other scheme is refused). Four things come with it:
\begin{itemize}
\item \textbf{The orbital energies, rebuilt at every point.} block2 runs
  \code{canonicalized = True} and uses the \code{mo\_energy} it is handed, so
  every evaluation hands it Eq.~\eqref{eq:fock}'s diagonal at that $(C, c)$.
\item \textbf{The 4-RDM, rebuilt when $c$ moves.} It is cached across the
  $\kappa$ sweep, keyed on the CI vector's bytes.
\item \textbf{A PC answer must exist.} \code{\_check\_pc\_conditioning} refuses
  LiF's state as a continuum through the cut, so \code{--pc} on LiF stops at
  the reference state. There is no PC gradient check on LiF, by design
  (\code{docs/testing.md}, ``LiF has no PC gradient, and that is the answer'').
\item \textbf{Its own CI step, and the guard at every displaced CI vector.} A
  determinant-axis displacement is not a spin eigenfunction. It lifts the PC
  metric's structurally null modes, and a lifted mode past the cut is kept and
  divided by. That is an error of order $h/h$ in the quotient, which no step
  ladder converges away and no residual sees. \code{CI\_STEP\_PC} keeps the
  lifted modes below the cut. Every PC energy at a CI vector other than the
  state's own runs the conditioning guard, and a rank test against the
  reference state's smallest kept singular value
  (\code{PC\_RANK\_KEPT\_FRACTION}) catches a mode kept outright. A ledger
  counts evaluations and guard runs in two different places, so a path that
  skips the guard fails the run (\code{check\_pc\_guard\_ledger}).
\end{itemize}

\begin{algorithm}[htbp]
\caption{Path 2, PC on CASCI orbitals}
\label{alg:path2}
\begin{algorithmic}[1]
\Statex \textit{Offline}: Algorithm~\ref{alg:path1}'s lines 1--2, \code{nevpt2\_numgrad.py --pc --npz \textit{N}} (no sidecar is written)
\Statex \textit{GPU} (\code{nevpt2\_sc\_pseudodensity ... --pc --dump-source-terms \textit{P}}): Algorithm~\ref{alg:dump} with PC's walk
\Statex \textit{Offline} (\code{nevpt2\_grad.py --case \textit{c} --pc [--source dump|both --source-path \textit{P}] --grad-num \textit{N}})
\State Algorithm~\ref{alg:path1}'s offline lines, with block2's $\Ecorr$, the PC steps of Table~\ref{tab:steps}, and the guard at every displaced $c$
\State \textbf{PASS} iff Eq.~\eqref{eq:gate} holds against \code{grad\_num\_pc}, every residual-independent check passes, and the guard ledger balances
\end{algorithmic}
\end{algorithm}

The equations are Path~1's, with Eqs.~\eqref{eq:pcdot} and \eqref{eq:pcadj}
in place of Eq.~\eqref{eq:pqr} on the GPU side. The cases are acrolein
CAS(6,5) and N2 CAS(4,4). The cost is the $\kappa$ sweep, one block2 PC solve
per non-gauge rotation (\code{--workers N} forks it, bit-identically to the
serial loop). The verdicts are \code{docs/testing.md}, ``The analytic gradient
against the numerical one'', and \code{docs/pc-nevpt2.md}, ``The source terms
by reverse mode''.

%======================================================================
\section{Path 3: SC-NEVPT2 on a CASSCF reference}
\label{sec:sc-casscf}

\paragraph{Offline, before.} The golden is the energy's Path~3: a CASCI on the
committed \code{\_casscf.mogold} orbitals, which pass the \code{--mo-from}
stationarity gate (\code{docs/energy.tex}, Eq.~(39)). The CASSCF-reference
sidecar is written by \code{--mo-from --gradgold}. The numerical gradient is
\code{nevpt2\_numgrad.py --reference casscf}, which re-optimises every
displaced state (\S\ref{sec:numgrad}).

\begin{algorithm}[htbp]
\caption{Path 3, SC on CASSCF orbitals}
\label{alg:path3}
\begin{algorithmic}[1]
\Statex \textit{Offline} (\code{generate\_golden.py --mo-from \textit{m} --gradgold}; \code{nevpt2\_numgrad.py --reference casscf --npz \textit{N}})
\State golden and sidecar on the CASSCF orbitals
\State \code{grad\_num} $\gets$ Eq.~\eqref{eq:numgrad} over re-optimised, continuity-guarded displaced states
\Statex \textit{GPU}: Algorithm~\ref{alg:dump} on the \code{\_casscf} golden, \textbf{unchanged}
\Statex \textit{Offline} (\code{nevpt2\_grad.py --case \textit{c} --reference casscf [--source both --source-path \textit{P}] [--grad-num \textit{N}]})
\State rebuild the state through \code{--mo-from}'s refusals (fingerprint, symmetry purity, the stationarity gate); assert \cons{O2}
\State read the dump unrotated \Comment{the state is the golden's basis}
\State $(z^c, Z, z) \gets$ Algorithm~\ref{alg:casscf};\ $g \gets$ Algorithm~\ref{alg:kernel}
\State the residual-independent checks, plus the Hessian-vector ladder \Comment{\S\ref{sec:checks}}
\State \textbf{PASS} iff Eq.~\eqref{eq:gate} holds (where a \code{grad\_num} exists) and every check passes
\end{algorithmic}
\end{algorithm}

The equations are Eqs.~\eqref{eq:pqr}--\eqref{eq:dphi},
\eqref{eq:Lcasscf}--\eqref{eq:cp} and \eqref{eq:numgrad}--\eqref{eq:gate}. The
GPU side is reference-agnostic: source terms are derivatives at fixed
orbitals, and the dump takes a \code{casscf} golden behind the source-term gate.
The cases are N2 CAS(4,4) (the inactive-degeneracy probe), LiF CAS(4,6), and
HSOH CAS(6,6), the $C_1$ case. HSOH keeps its point group under every
displacement, so all $3N$ components are probed and none is zero by symmetry.
Its committed \code{\_casscf.gradgold} carries the re-optimised sweep's
\code{grad\_num}, which is read by default. Without a \code{grad\_num} the gate
is the residual-independent checks. The verdicts are \code{docs/testing.md},
``SC on a CASSCF reference: LiF, and the N2 degeneracy probe'' and ``SC on a
CASSCF reference: the C1 case''.

%======================================================================
\section{Path 4: PC-NEVPT2 on a CASSCF reference}
\label{sec:pc-casscf}

\paragraph{What it is.} Path~2's \code{--pc} fed into Path~3's coupled solve.
$Y$, $y$ and $\Phi$'s $\Ecorr$ are the PC energy's, at the PC steps and under
the guard. $J$, \cons{S1}, \cons{C1}, \cons{O2} and $B$ are the reference's
and do not change.

\begin{algorithm}[htbp]
\caption{Path 4, PC on CASSCF orbitals}
\label{alg:path4}
\begin{algorithmic}[1]
\Statex \textit{Offline}: Algorithm~\ref{alg:path3}'s sweep; then \code{pc\_on\_states.py} (the GPU's PC energy on each kept state) and \code{nevpt2\_numgrad.py --reference casscf --pc-states \textit{S} --npz \textit{N}}
\Statex \textit{GPU}: Algorithm~\ref{alg:dump} with \code{--pc}, on the \code{\_casscf} golden
\Statex \textit{Offline} (\code{nevpt2\_grad.py --case \textit{c} --reference casscf --pc [--source both --source-path \textit{P}] --grad-num \textit{N}})
\State Algorithm~\ref{alg:path3}'s offline lines, with Path~2's changes
\State \textbf{PASS} iff Eq.~\eqref{eq:gate} holds against \code{grad\_num\_pc}, every check passes, and the guard ledger balances
\end{algorithmic}
\end{algorithm}

The equations are Path~3's with Path~2's substitutions. Its numerical gradient
is the GPU's PC energy, differenced over the same re-optimised states the SC
sweep kept, so the oracle's PC energies come from the card while its states
come from PySCF. The case is acrolein CAS(6,5), block2-sourced and
GPU-sourced (\code{docs/testing.md}, ``PC on a CASSCF reference: acrolein'').

%======================================================================
\section{Checks that read no residual}
\label{sec:checks}

A converged residual says that a solve found the answer to the equations it
was given, not that those were the right equations. Every solve and every
difference quotient in \code{nevpt2\_grad.py} therefore also gets a check that
reads no residual. Each such check runs through \code{Functional.evaluate},
the route $\Phi$ is differentiated through, rather than through the code under
test. Most are \emph{step ladders}: central differences at $h = 2^{-k}$ over a
range of $k$, gated on two facts.
\begin{enumerate}
\item \textbf{The ladder converges.} The change between neighbouring rungs falls
  by a factor in $(3.2, 4.8)$ for at least three consecutive halvings, and of
  several such runs the finest is used. This is the C++ checks' rule too,
  \code{nevpt2::longestFourfoldRun}.
\item \textbf{It converges to the right number.} The Richardson extrapolant at
  that run's finest pair is within \code{LADDER\_TOL} $= 10^{-6}$ of the
  quantity under test. Comparing an extrapolant, rather than the smallest
  $|f_d - \mathrm{ref}|$ in the sweep, is what keeps a crossing of a wrong
  reference from reading as agreement.
\end{enumerate}

\begin{table}[htbp]
\centering
\small
\begin{tabular}{@{}p{0.24\textwidth}p{0.44\textwidth}p{0.26\textwidth}@{}}
\toprule
check & ladder of & converges to \\
\midrule
projection of $y$, of $z$ & --- & $|c\cdot y|$, $|c\cdot z| \le 10^{-6}$ \\
$y$ linearity (two seeds) & $\Ecorr$ along a spin-adapted, in-sector $u\perp c$, on the sphere & $\langle u, y\rangle$ \\
CI stationarity & Eq.~\eqref{eq:Lcasci}'s CI part along $u$ (whole, or by pieces) & $0$ \\
orbital stationarity & $L$ along a totally symmetric non-gauge $\kappa$, $C e^{t\kappa}$ & $0$ \\
nuclear ladder & $\Phi$ along a totally symmetric displacement & $\langle d, g\rangle$ \\
PC guard ledger & --- & every displaced $c$ guarded, no rank change \\
Hessian-vector (CASSCF) & \code{gen\_g\_hop}'s $g$ along seeded random $v$ & $Jv$ \\
\bottomrule
\end{tabular}
\caption{The residual-independent checks (\code{check\_*} in
\code{nevpt2\_grad.py}). The non-gauge rotations of the orbital-stationarity
check are derived from the RHF occupied projector, independently of the
solve's \code{kappa\_pairs}, so the check does not inherit that bookkeeping.
Which check catches which of the errors found so far, with the numbers, is
\code{docs/testing.md}, ``Residual-independent checks of the gradient
path''.}
\label{tab:checks}
\end{table}

\paragraph{Validation with $\Ecorr$ dropped.} \code{--casci-only} runs the
identical pipeline with $\Ecorr$ removed from the source terms and from $\Phi$.
On the CASCI branch the reference is then a plain CASCI, and its gradient is
gated against \code{nevpt2\_numgrad.py}'s \code{grad\_casci}. That one
comparison exercises the frozen path, the full constraint Jacobian, the
selection-boundary multipliers and the CPHF block. It does not exercise the
CI Z-vector (a CASCI is variational in $c$, so $y = z = 0$) or Eqs.~56--57,
whose numerators vanish for $\Ecas$. \code{pyscf.grad.casci} is an oracle only
for an energy-ordered active space (\code{--energy-ordered}). On the CASSCF
branch the reference becomes a CASSCF, stationary in $\kappa$ and $c$, so the
multipliers must come out \emph{exactly} zero and the gradient must be
\code{pyscf.grad.casscf}'s (\code{check\_multipliers\_vanish}).

\paragraph{The gauge identity} (\code{docs/gradient-theory.md}, \S1.3) is
reported, not gated: the Lagrangian's derivative along the gauge rotations.
On the CASCI branch the constraints' own gauge response is not part of the
sum, because $f$'s active sum is not covariant at fixed $\G1$. On the CASSCF
branch the constraints are covariant and the identity is the whole
Lagrangian's.

%======================================================================
\section{Scope and limits}
\label{sec:limits}

\begin{itemize}
\item \textbf{No nuclear derivative in \code{src/} or \code{apps/}.} The GPU
  computes and writes the source terms. The Z-vector solve, the skeleton
  derivative and the numerical gradient are offline Python. The in-tree
  Lagrangian gate, \code{requireLagrangianReference}, refuses a \code{casscf}
  golden, and no path in \code{src/} sits behind it yet
  (\code{docs/gradients.md}, \S1).
\item \textbf{The skeleton derivative is differenced, not contracted.} An AO
  contraction against derivative integrals would need the NEVPT2 relaxed
  densities: $Y$ fixes only the antisymmetric part of the generalised Fock.
  \code{src/gradient/} builds the pseudodensities and frees them as the slab
  walk ends, and no flag writes them out.
\item \textbf{The CASCI branch's selection is a discrete function of geometry.}
  Where it changes, the energy has a kink. A small gap across an active-space
  boundary makes the selection multipliers large. Validation therefore uses
  in-plane displacements of planar $C_s$ molecules, or a $C_1$ one on the
  CASSCF branch (\code{docs/gradient-theory.md}, \S4). On the CASSCF branch
  the analogue is the conditioning of $J$, and a displaced optimisation that
  lands on another branch.
\item \textbf{SC's non-invariance among degenerate inactive orbitals} (Park)
  reaches the gradient through the \cons{O2} divisions on both branches. N2 is
  the deliberate probe of it, and degenerate pairs follow the rule of
  \S\ref{sec:casci}.
\item \textbf{PC can refuse.} PC refuses a state whose metric is not gapped,
  with LiF as the standing example, and a displaced CI vector that lifts a null
  mode past the cut.
\item \textbf{Density fitting is SC, CASCI branch, one case.} There is no
  fitted PC dump and no fitted CASSCF \code{grad\_num}.
\item \textbf{The second derivative is elsewhere.} The energy's parameter
  Hessian is in tree; the nuclear Hessian is not. See
  \code{docs/hessian-theory.md} and \code{docs/gradients.md}, \S1.1.
\item \textbf{Cost.} The offline cost is the $\kappa$ sweep, $O(n_{\mathrm{mo}}^2)$
  energy evaluations, and that is the part the GPU's dump replaces for $Y$.
  The walls are \code{docs/reference-data.md}, ``The analytic gradient'', and
  \code{docs/testing.md}, ``PC on a CASSCF reference: acrolein''.
\end{itemize}

%======================================================================
\appendix
\section{The source-term inventory}
\label{app:inventory}

$Y$ reads, besides what the energy reads, the arrays the promotion rule
generates: every stored block with one index replaced by a general one $x$,
keeping the charge-distribution pairing (\code{docs/gradient-theory.md},
\S3.1--3.3). Write $(xP|QR)$ for the array in which $x$ shares a charge
distribution with a label from space $P$. Seventeen of the eighteen are
needed, and the sidecar carries them:

\begin{table}[htbp]
\centering
\small
\begin{tabular}{@{}lll@{}}
\toprule
array & shape & differentiates \\
\midrule
$(xa|aa)$ & $n_{\mathrm{mo}}n_a^3$ & $aaaa$, $aaav$, $aaac$ \\
$(xv|aa)$ & $n_{\mathrm{mo}}n_vn_a^2$ & $aaav$, $caav$, $\eps_r$ \\
$(xa|va)$ & $n_{\mathrm{mo}}n_a^2n_v$ & $aaav$, $aavv$, $caav$, $\eps$ \\
$(xc|aa)$ & $n_{\mathrm{mo}}n_cn_a^2$ & $aaac$, $caav$, $\heff$, $\eps_i$ \\
$(xa|ca)$ & $n_{\mathrm{mo}}n_a^2n_c$ & $aaac$, $ccaa$, $caav$, $\eps$ \\
$(xv|va)$ & $n_{\mathrm{mo}}n_v^2n_a$ & $aavv$, $cavv$ \\
$(xc|ca)$ & $n_{\mathrm{mo}}n_c^2n_a$ & $ccaa$, $ccav$, $\heff$ \\
$(xa|cv)$ & $n_{\mathrm{mo}}n_an_cn_v$ & $caav$, $ccav$, $cavv$ \\
$(xv|ca)$ & $n_{\mathrm{mo}}n_vn_cn_a$ & $caav$, $ccav$ \\
$(xc|va)$ & $n_{\mathrm{mo}}n_cn_vn_a$ & $caav$, $cavv$, $\heff$ \\
$(xc|cv)$ & $n_{\mathrm{mo}}n_c^2n_v$ & $ccav$, $ccvv$, $\heff$ \\
$(xv|cv)$ & $n_{\mathrm{mo}}n_v^2n_c$ & $cavv$, $ccvv$ \\
$(xa|cc)$ & $n_{\mathrm{mo}}n_an_c^2$ & $\heff$ (exchange), $\eps_i$ \\
$(xv|cc)$ & $n_{\mathrm{mo}}n_vn_c^2$ & $\heff$ (exchange), $\eps_r$ \\
$(xc|cc)$, two core labels tied & $n_{\mathrm{mo}}n_c^2$ & $\heff$, $\eps_i$ \\
$(xa|vv)$, virtual labels tied & $n_{\mathrm{mo}}n_an_v$ & $\eps_r$ \\
$(xc|vv)$, virtual labels tied & $n_{\mathrm{mo}}n_cn_v$ & $\eps_r$ \\
\bottomrule
\end{tabular}
\caption{The one-general-index arrays (\code{docs/gradient-theory.md}, \S3.3).
The eighteenth, $(xv|vv)$, is never needed: no integral the energy reads has
three virtual labels, and promotion cannot add one.}
\label{tab:inventory}
\end{table}

Besides the integrals, the sidecar carries the full-range $h$, $\heff$ and its
variants, $f$, $f^{\mathrm h}$, \code{mo\_coeff} and, for the CASCI branch
only, the RHF Fock $F$ and the RHF electron count
(\code{docs/gradient-theory.md}, \S3.5). Three invariants are checked on it.
$F$ is block-diagonal over core/active/virtual (\cons{O1}; CASCI only). $f$'s
core and virtual off-diagonals vanish, and its diagonal is \code{e\_core},
\code{e\_virt} (\cons{O2}). $F$'s active block splits into the right number of
occupied and virtual eigenvalues (CASCI only). Density-fitted, each array is
$\sum_L B^L_{xp}B^L_{qr}$ (Eq.~\eqref{eq:dB}). It reads all six space blocks
of $B$, and the \code{NEVPT2F1} sidecar adds the two that a \code{--df} golden
lacks, $B_{cc}$ and $B_{vv}$ (\code{docs/gradient-theory.md}, \S3.8).

\section{\texorpdfstring{The pseudocanonical multipliers $z^c$}{The pseudocanonical multipliers z\^{}c}}
\label{app:z}

The energy reads $f$ only on its core and virtual diagonals, so $f$'s
off-diagonal response there must be pinned by the convention \cons{O2}. In the
pseudocanonical gauge $f$ is diagonal on both spaces, so for a rotation inside
one of them Eq.~\eqref{eq:dh} collapses to $\delta f_{pq} =
\kappa_{qp}(\eps_q - \eps_p)$. The \cons{O2} rows of the stationarity
equations then reduce to one division per pair, Park's Eqs.~56--57:
\begin{equation}
  z^c_{ij} = \frac12\,\frac{Y_{ij} - Y_{ji}}{\eps_i - \eps_j}\quad(i\ne j\ \text{core}),\qquad
  z^c_{rs} = \frac12\,\frac{Y_{rs} - Y_{sr}}{\eps_r - \eps_s}\quad(r\ne s\ \text{virtual}).
  \label{eq:z5657}
\end{equation}
The code does not divide by hand. It solves \cons{O2}'s rows of $M$
(Eqs.~\eqref{eq:orbz} and \eqref{eq:zc}), whose diagonal holds these
denominators, and the degenerate-pair rule of \S\ref{sec:casci} replaces a
division the state does not determine. Neither branch adds anything to this
appendix, which is why \code{docs/gradient-theory.md}, \S2.3 calls it shared.
The GPU divides the same two blocks of its own $Y$ in
\code{nevpt2\_sc\_pseudodensity --kappa-full}, where it is a check on $Y$
rather than an input to a gradient (\code{docs/gradients.md}, \S6).

\end{document}
